\documentclass{article}
\usepackage[utf8]{inputenc}
\usepackage[T1]{fontenc}   % fixes broken accents in the PDF text layer (Janó Muñoz, Zúñiga-Fernández, Autònoma)
\usepackage{lmodern}
\usepackage[margin=1.5cm]{geometry}
\usepackage{natbib}
\usepackage{amssymb, amsmath}
\usepackage{xcolor}

\usepackage[colorlinks=True,urlcolor=blue]{hyperref}
\usepackage{array} % for defining new column types
\usepackage{fontawesome} % for icons
\usepackage{comment} % for commenting out sections
\usepackage{titlesec} % for custom section titles

% Table formatting
\usepackage{longtable}
\setlength\LTleft{0pt}      % left-align (longtable centres by default)
\setlength\LTright{\fill}
\setlength\LTpre{0pt}       % remove the \bigskip longtable adds above...
\setlength\LTpost{0pt}      % ...and below

% Define new column types for the tables
\newcolumntype{L}{>{\raggedleft}p{0.175\textwidth}}
\newcolumntype{R}{p{0.765\textwidth}}

% Custom section titles
\titleformat{\section}{\Large\bfseries}{\thesection}{}{}[{\titlerule[0.5pt]}]
\titleformat{\subsection}{\normalsize\bfseries}{\thesection}{}{}
\titlespacing*{\subsection}{0pt}{1.2ex}{0.6ex}

% Author-position tag for co-authored papers
\newcommand{\me}{\textbf{Janó Muñoz, C.}}
\newcommand{\pos}[2]{\,(\me; #1 of #2 authors)}

\begin{document}

\small % Reduce font size for the whole document

\begin{center}
\huge
\textbf{Clàudia Janó Muñoz} \\
\normalsize
PhD student - University of Cambridge \\
\end{center}

\begin{minipage}[ht]{0.6\linewidth}
	\faEnvelope~~\textbf{Email}: cj467@cam.ac.uk / claudiajano0@gmail.com\\
	\faGithub~~\textbf{Github:} \href{https://github.com/claudiajanomunoz}{github.com/claudiajanomunoz}\\
	\faGlobe~~\textbf{Website:} \href{https://claudiajanomunoz.github.io}{claudiajanomunoz.github.io} \\
  \faPhone~~\textbf{Mobile:} +44(0)7935493960 \\
  \faMapMarker~~\textbf{Address:} 19, J J Thomson Avenue, Cambridge, CB3 0FA, UK \\
  \textbf{Date of Birth:} 30/11/2000 \\
  \textbf{Nationality:} Spanish \\
\end{minipage}
\begin{minipage}[ht]{0.34\linewidth}
	\begin{flushright}
    \textbf{ORCID:} \href{https://orcid.org/0009-0008-5713-0750}{0009-0008-5713-0750} \\
    \textbf{arXiv:} \href{http://arxiv.org/a/janomunoz_c_1}{janomunoz\_c\_1} \\
    \textbf{Google Scholar:} \href{https://scholar.google.com/citations?user=7sh8LMYAAAAJ&hl=en&oi=sra}{Clàudia Janó Muñoz} \\
    \textbf{NASA/ADS:} \href{https://ui.adsabs.harvard.edu/search/q=author%3A%22Janó%20Muñoz%2C%20C.%22&sort=date%20desc%2C%20bibcode%20desc&p_=0}{Clàudia Janó Muñoz} \\
	\end{flushright}
\end{minipage}

\vspace{0.6ex}
\noindent\textbf{UAT keywords:} Transit photometry; Astronomical detectors; M dwarf stars % CHECK: confirm exact terms at astrothesaurus.org (max. 3; required by the Clay call)

\section*{Research interests}
% DRAFT - edit freely. Keep to 4-5 lines; end on the direction of your Clay proposal.
\noindent I find and characterise transiting exoplanets around M dwarfs and ultracool dwarfs with ground-based near-infrared photometry. My PhD centres on SPIRIT, an InGaAs CMOS infrared camera at the SPECULOOS Southern Observatory: I optimised its calibration, characterised its on-sky noise budget and performance, and used it to validate new planets such as the warm Neptune TOI-2407\,b. I am interested in pushing ground-based photometric precision through detector-level understanding of instrumental systematics, and in using that precision to build complete samples of small planets around the nearest low-mass stars. % CHECK: final sentence should match your Clay research proposal

\section*{Education}
\begin{tabular}{L!{\vrule}R}
	2022 -- 2027 (exp.) & \textbf{PhD in Physics}, Cavendish Laboratory, University of Cambridge. Thesis: \textit{Optimisation, performance and first discoveries of the SPIRIT infrared detector for exoplanets transiting late-type stars}. Supervised by Prof. Didier Queloz and Dr. Matthew Hooton. Expected thesis submission: February 2027. \\ % CHECK: submission vs. viva month
	2018 -- 2022 & \textbf{MPhys Physics with Astrophysics}, First Class, University of Manchester.  \\
\end{tabular}

\section*{Research experience}
\begin{tabular}{L!{\vrule}R}
    Oct. 2022 -- present & \textbf{PhD research}, Cavendish Laboratory, University of Cambridge (SPECULOOS consortium). \newline
    \hspace*{0.6em}$\bullet$ Optimised the calibration and characterised the on-sky performance of SPIRIT, an InGaAs CMOS near-infrared camera on the Callisto telescope (SPECULOOS-South): dark current and glow, row-pattern correction, bad-pixel mapping, noise model, and the $-60$ vs.\ $-50\,^{\circ}$C operating trade-off. \newline
    \hspace*{0.6em}$\bullet$ Analysed the first year of SPIRIT operations: 328 visits of 68 targets (ultracool dwarfs and TESS candidates). \newline % CHECK: counts must match the final SPIRIT paper
    \hspace*{0.6em}$\bullet$ Validated the warm Neptune TOI-2407\,b with multi-band ground-based follow-up including SPIRIT (first-author, MNRAS 2025). \newline
    \hspace*{0.6em}$\bullet$ Co-developing \texttt{goibniu}, an observation-scheduling tool (Hooton \& Janó Muñoz, in prep.). \\ % CHECK: wording of your role
    Oct. 2021 -- May 2022 & \textbf{Master's project}, University of Manchester. Machine learning in spectral analysis of ALMA observations. Supervised by Prof. Gary Fuller.\\
    Sep. 2021 -- Oct. 2021 & \textbf{Summer research project}, Universitat Autònoma de Barcelona. Simulating the effects of close stellar encounters with N-body simulators. Supervised by Dr. Ignasi Ribas. \\
    Jun. 2021 -- Aug. 2021 & \textbf{UROP Summer Internship}, Imperial College London. Exoplanet atmospheric evolution with machine learning (led to Rogers, Janó Muñoz et al. 2023, MNRAS). Supervised by Dr. James Owen.\\
\end{tabular}

\section*{Grants and awards}
\begin{tabular}{L!{\vrule}R}
    October 2022 & \textbf{Harding Distinguished Postgraduate Scholarship} (£20,000 per year of study), including a £2,000 yearly stipend for academic travel. \\
    December 2018 & \textbf{Physics Success Scholarship} (£2,000), for academic excellence in physics and maths.\\
\end{tabular}

\section*{Observing and proposals}
\begin{tabular}{L!{\vrule}R}
  July 2026 & \textbf{Co-PI}, \textit{Unlocking Ultracool Dwarf Magnetospheres via Simultaneous Radio/IR Monitoring}, MeerKAT radio telescope. \\ % CHECK: add outcome (pending / awarded N h)
  Oct. 2022 -- Jan. 2025 & Regular remote observing with the Callisto and Ganymede telescopes, SPECULOOS Southern Observatory. \\ % CHECK: replace with a number of nights (original said "bi-monthly", which is ambiguous)
  Dec. 2023 & On-site observations, instrument maintenance and testing at the SPECULOOS Southern Observatory, Cerro Paranal, Chile. \\
\end{tabular}

\section*{Services}
\subsection*{Committees}
\begin{tabular}{L!{\vrule}R}
  Jan. 2025 -- Apr. 2025 & \textbf{International Women's Day Committee}, Member, Institute of Astronomy, Cambridge. \\
\end{tabular}
% CHECK: add refereeing, journal-club or LOC roles here if you have any.

\subsection*{Outreach}
\begin{tabular}{L!{\vrule}R}
  Mar. 2025 & \textbf{International Women's Day}, Ran an exoplanets workshop for high-school students, Institute of Astronomy, Cambridge.\\
  Mar. 2025 & \textbf{Open Day}, Volunteer, Institute of Astronomy, Cambridge.\\
\end{tabular}

\section*{Teaching responsibilities}
\begin{tabular}{L!{\vrule}R}
	Jun. 2026 -- Sep. 2026 & \textbf{Teaching Assistant}, Cambridge Centre for International Research, UK. \\
	Apr. 2024 -- May 2024 & \textbf{Part III Example Classes: Exoplanets Course}, Cavendish Laboratory, University of Cambridge. \\
  Nov. 2022 -- May 2023 & \textbf{Part IA Physics Supervisions}, Emmanuel College, University of Cambridge.\\
\end{tabular}

\section*{Talks}
\subsection*{Seminars} % CHECK: rename to "Invited talks" if these were invitations
\begin{tabular}{L!{\vrule}R}
    March 2026 & \textit{SPIRIT: A CMOS Detector for Exoplanets Around Ultracool Dwarfs}, Center for Astrophysics $|$ Harvard \& Smithsonian, Cambridge, MA, USA. \\
    November 2025 & \textit{SPIRIT: A CMOS Detector for Exoplanets Around Ultracool Dwarfs}, Exoplanet Seminar, Observatoire de Paris-PSL, France. \\
    August 2022 & \textit{Emulating photoevaporation with Machine Learning}, University of California Santa Cruz, USA (remote). \\
\end{tabular}

\subsection*{Conference talks}
\begin{tabular}{L!{\vrule}R}
    February 2026 & \textit{Search for habitable exoplanets with the SPIRIT infrared detector}, Harding Scholar Conference, Cambridge, UK. \\
    January 2026 & \textit{First Results from SPIRIT: A CMOS Detector for Exoplanets around Ultracool Dwarfs}, Rocky Worlds IV, Groningen, Netherlands. \\
\end{tabular}

\subsection*{Posters}
\begin{tabular}{L!{\vrule}R}
    June 2026 & \textit{Ground-based transits in the infrared with SPIRIT}, Exoplanets VI, Porto, Portugal. \\
    June 2024 & \textit{Ground-based transits in the infrared with SPIRIT}, Exoplanets V, Leiden, Netherlands. \\
    January 2024 & \textit{Finding transiting planets in the infrared with SPIRIT}, Rocky Worlds III, ETH Zurich, Switzerland. \\
\end{tabular}

\subsection*{Departmental talks}
\begin{tabular}{L!{\vrule}R}
    December 2025 & \textit{Discussion on SPIRIT: Infrared CMOS Photometry for Exoplanets around Ultracool Dwarfs}, LCLU discussion series, Cambridge, UK. \\
    November 2025 & \textit{Discovery of the Neptune-like planet TOI-2407\,b}, Harding Scholars Internal Speaker Event, St Catharine's College, Cambridge, UK. \\
    February 2025 & \textit{Exoplanet Detection with SPIRIT: Infrared CMOS Photometry and the Discovery of the Hot Neptune TOI-2407\,b}, Hills Coffee Talks, Battcock Centre for Experimental Astrophysics, Cambridge, UK. \\
    February 2025 & \textit{TOI-2407\,b: A hot Neptune in the Neptune Desert discovered with the SPIRIT Infrared Detector}, Exoplanet Seminar series, University of Cambridge, UK. \\
\end{tabular}

\section*{Software development}
\begin{tabular}{L!{\vrule}R}
  \texttt{goibniu} & Co-developer: observation-scheduling tool for SPECULOOS (Hooton \& Janó Muñoz, in prep.). \\ % CHECK: add link if public
  \href{https://github.com/claudiajanomunoz/SMART}{\texttt{SMART}} & Co-author: machine-learning tool for fitting protostellar spectra. \\
  % CHECK: add your SPIRIT calibration/reduction code if it is a distinct package, and your contribution to Astra (Pedersen et al. 2026) if you wrote part of it.
\end{tabular}

\section*{Extracurricular activities}
\begin{tabular}{L!{\vrule}R}
  May 2021 -- Feb. 2022 & \textbf{Astronomy Society}, Chair, University of Manchester. \\
  Feb. 2020 -- May 2021 & \textbf{Astronomy Society}, Science Officer, University of Manchester. \\
\end{tabular}

\section*{Skills}
\begin{tabular}{L!{\vrule}R}
	Programming & \textbf{Proficient}: Python; \textbf{Intermediate}: C++\\
	Computing & \textbf{Experienced}: LaTeX, Linux, Docker; \textbf{Basic}: HPC \\
	Languages & \textbf{Native}: Catalan, Spanish; \textbf{Fluent}: English; \textbf{Intermediate}: German (B1)
\end{tabular}

\section*{Publications}
\noindent 10 refereed or accepted journal papers (1 first-author, 1 second-author), 1 first-author paper under review, 2 conference proceedings. % CHECK: recount at submission

\subsection*{First and second author}
\begin{tabular}{L!{\vrule}R}
    under review &
    \me, et al.
    {\textit{Performance and prospects of the SPIRIT infrared CMOS detector for exoplanet detection}}, MNRAS \\ % CHECK: "submitted to MNRAS" + arXiv ID if submitted; otherwise "in prep."

    June 2025 &
    \me, et al.
    \href{https://ui.adsabs.harvard.edu/abs/2025MNRAS.541..630J/abstract}{\textit{TOI-2407 b: a warm Neptune in the desert}},
    Monthly Notices of the Royal Astronomical Society, 541, 630--641 \\

    March 2023 &
    Rogers, J. G., \me, Owen, J. E., Makinen, T. L.
    \href{https://ui.adsabs.harvard.edu/abs/2023MNRAS.519.6028R/abstract}{\textit{Exoplanet atmosphere evolution: emulation with neural networks}},
    Monthly Notices of the Royal Astronomical Society, 519, 6028--6043 \\
\end{tabular}

\subsection*{Contributing author}
% CHECK: optionally add your contribution after each entry, e.g. "SPIRIT observations and photometry".
\begin{longtable}{L!{\vrule}R}
    September 2026 &
    Scott, M. G., et al.\pos{33rd}{45}
    \href{https://ui.adsabs.harvard.edu/abs/2026arXiv260937597S/abstract}{\textit{TOI-6981 b: a transiting sub-Neptune at the water ice line}},
    accepted to Monthly Notices of the Royal Astronomical Society, arXiv:2609.37597 \\

    September 2026 &
    Lang, F. Z., et al.\pos{37th}{47}
    \href{https://ui.adsabs.harvard.edu/abs/2026MNRAS.tmp.1622L/abstract}{\textit{TOI-4616 b: an Earth-sized planet transiting a nearby M4 dwarf}},
    Monthly Notices of the Royal Astronomical Society \\ % CHECK: arXiv title reads "a benchmark Earth-sized planet"; use the published title

    April 2026 &
    Dransfield, G., et al.\pos{13th}{50}
    \href{https://ui.adsabs.harvard.edu/abs/2026MNRAS.547ag448D/abstract}{\textit{MANGOS -- II. Five new giant planets orbiting low-mass stars}},
    Monthly Notices of the Royal Astronomical Society, 547, stag448 \\

    March 2026 &
    Scott, M. G., et al.\pos{27th}{33}
    \href{https://ui.adsabs.harvard.edu/abs/2026MNRAS.547ag070S/abstract}{\textit{Two temperate Earth- and Neptune-sized planets orbiting fully convective M dwarfs}},
    Monthly Notices of the Royal Astronomical Society, 547, stag070 \\

    December 2025 &
    Barkaoui, K., et al.\pos{26th}{44}
    \href{https://ui.adsabs.harvard.edu/abs/2025MNRAS.544.2637B/abstract}{\textit{TOI-7166 b: a habitable zone mini-Neptune planet around a nearby low-mass star}},
    Monthly Notices of the Royal Astronomical Society, 544, 2637 \\

    October 2025 &
    Zúñiga-Fernández, S., et al.\pos{38th}{66}
    \href{https://ui.adsabs.harvard.edu/abs/2025A%26A...702A..85Z/abstract}{\textit{Two warm Earth-sized exoplanets and an Earth-sized candidate in the M5V-M6V binary system TOI-2267}},
    Astronomy \& Astrophysics, 702, A85 \\

    July 2024 &
    Gillon, M., et al.\pos{23rd}{84}
    \href{https://ui.adsabs.harvard.edu/abs/2024NatAs...8..865G/abstract}{\textit{Detection of an Earth-sized exoplanet orbiting the nearby ultracool dwarf star SPECULOOS-3}},
    Nature Astronomy, 8, 865--878 \\

    October 2023 &
    Triaud, A. H. M. J., et al.\pos{20th}{43}
    \href{https://ui.adsabs.harvard.edu/abs/2023MNRAS.525L..98T/abstract}{\textit{An M dwarf accompanied by a close-in giant orbiter with SPECULOOS}},
    Monthly Notices of the Royal Astronomical Society: Letters, 525, L98--L104 \\

    July 2026 &
    Pedersen, P. P., et al.\pos{11th}{15}
    \href{https://ui.adsabs.harvard.edu/abs/2026arXiv260712898P/abstract}{\textit{Astra: an open-source fully autonomous robotic observatory control software}},
    arXiv e-prints, arXiv:2607.12898 \\ % CHECK: add "submitted to <journal>" if applicable
\end{longtable}

\subsection*{Conference proceedings}
\begin{tabular}{L!{\vrule}R}
    August 2024 &
    Zúñiga-Fernández, S., et al.\ (incl.\ \me)
    \href{https://ui.adsabs.harvard.edu/abs/2024SPIE13094E..13Z/abstract}{\textit{SPECULOOS: five years hunting terrestrial planets around ultra-cool dwarfs}},
    Proceedings of the SPIE, 13094 \\ % CHECK: author position (not on arXiv)

    July 2024 &
    Pedersen, P. P., et al.\pos{11th}{18}
    \href{https://ui.adsabs.harvard.edu/abs/2024SPIE13096E..3XP/abstract}{\textit{Infrared photometry with InGaAs detectors: first light with SPECULOOS}},
    Proceedings of the SPIE, 13096 \\
\end{tabular}

\end{document}
